%%%%%%%%%%  scan idx 435  |  folio 422  %%%%%%%%%%
\SGAPrintPage{422}%
\sgaStableTarget{sga6.e07.prop.1-7}\textbf{Proposition 1.7.} (\uline{Compare} EGA IV 19.1.5 (iv)). \uline{Let} $i:Y\longrightarrow X$, \uline{and} $j:Z\longrightarrow Y$ \uline{be two immersions}:
\begin{enumerate}[label=\roman*)]
\item \uline{If} $i$ \uline{and} $j$ \uline{are regular}, $i\circ j$ \uline{is regular}, \uline{and the sequence}:
\[
0\longrightarrow j^*(\underline{N}_{X/Y})\longrightarrow \underline{N}_{X/Z}\longrightarrow \underline{N}_{Y/Z}\longrightarrow0
\]
\uline{is exact}; \uline{in particular}, $\codim(i\circ j)=\codim(i)+\codim(j)$.
\item \uline{If} $i\circ j$ \uline{and} $i$ \uline{are regular}, \uline{and if the sequence}:
\[
0\longrightarrow j^*(\underline{N}_{X/Y})\longrightarrow \underline{N}_{X/Z}\longrightarrow \underline{N}_{Y/Z}\longrightarrow0,
\]
\uline{is exact and locally split}, $j$ \uline{is regular};
\item \uline{If} $X$ \uline{is noetherian}, \uline{and if} $i\circ j$ \uline{and} $j$ \uline{are regular}, $i$ \uline{is regular at the points of} $Z$.
\end{enumerate}

Since the question is local on $X$, we may suppose $X$ affine and the immersions closed; let $X=\Spec(A)$, $Y=\Spec(B)$, $Z=\Spec(C)$, with $B=A/I$, $C=A/J=B/K$.

i) After localizing, we may suppose $I/I^2$ free over $B$, and $K/K^2$ free over $C$ (1.3). Let $f_1,\ldots,f_n$ be elements of $A$ lifting a basis of $I/I^2$, $g_1,\ldots,g_m$ elements of $A$ lifting a basis of $K/K^2$; then, after further localization, we may suppose that $f_1,\ldots,f_n$ generate $I$, that $f_1,\ldots,f_n,g_1,\ldots,g_m$ generate $J$, and that $\bar g_1,\ldots,\bar g_m$ generate $K$, where $\bar g_j$ is the image of $g_j$ in $B$ (1.4); moreover, $f_1,\ldots,f_n$ form a minimum system of generators of $I$, and $\bar g_1,\ldots,\bar g_m$ a minimum system of generators of $K$.

Consider the diagram:
\sgaStableTarget{sga6.e07.eq.1-7-1}\begin{equation}\tag{1.7.1}
\begin{tikzcd}[column sep=large,row sep=large]
A_0=\mathbb Z[T_1,\ldots,T_{n+m}] \arrow[r] \arrow[d] & A \arrow[d]\\
B_0=\mathbb Z[T_1,\ldots,T_m] \arrow[r] \arrow[d] & B \arrow[d]\\
C_0=\mathbb Z \arrow[r] & C,
\end{tikzcd}
\end{equation}

%%%%%%%%%%  scan idx 436  |  folio 423  %%%%%%%%%%
\SGAPrintPage{423}%
where the horizontal arrows are defined by sending $T_i$ to $f_i$ for $1\leq i\leq n$, and $T_{n+j}$ to $g_j$ for $1\leq j\leq m$, where the homomorphism $A_0\longrightarrow B_0$ is defined by sending $T_i$ to $0$ for $1\leq i\leq n$, and $T_{n+j}$ to $T_j$ for $1\leq j\leq m$, and the homomorphism $B_0\longrightarrow C_0$ by setting the $T_j$, $1\leq j\leq m$, to zero. We then have:
\[
 B=A\otimes_{A_0}B_0,
 \qquad
 C=A\otimes_{A_0}C_0=B\otimes_{B_0}C_0.
\]
The situation under consideration is therefore deduced from the situation $(A_0,B_0,C_0)$ by the base change $A_0\longrightarrow A$. The corresponding immersions $i_0$, $j_0$ and $i_0\circ j_0$ are regular; consequently, the homology groups of the Koszul complexes associated with the sequences $(f_1,\ldots,f_n)$, $(f_1,\ldots,f_n,g_1,\ldots,g_m)$ and $(\bar g_1,\ldots,\bar g_m)$ are respectively the $\Tor_i^{A_0}(B_0,A)$, $\Tor_i^{A_0}(C_0,A)$, $\Tor_i^{B_0}(C_0,B)$. Consequently:
\sgaStableTarget{sga6.e07.eq.1-7-2}\begin{equation}\tag{1.7.2}
\begin{aligned}
 i\ \text{regular} &\Longleftrightarrow \forall i\geq1,
 \Tor_i^{A_0}(B_0,A)=0;\\
 i\circ j\ \text{regular} &\Longleftrightarrow \forall i\geq1,
 \Tor_i^{A_0}(C_0,A)=0;\\
 j\ \text{regular} &\Longleftrightarrow \forall i\geq1,
 \Tor_i^{B_0}(C_0,B)=0.
\end{aligned}
\end{equation}

We then use the transitivity spectral sequence ([1], XVI 5)
\sgaStableTarget{sga6.e07.eq.1-7-3}\begin{equation}\tag{1.7.3}
E^2_{pq}=\Tor_p^{B_0}\bigl(C_0,\Tor_q^{A_0}(B_0,A)\bigr)
\Longrightarrow \Tor_n^{A_0}(C_0,A)..
\end{equation}

By the hypotheses of i), $E^2_{pq}=0$ for $p$ or $q\ne0$. It therefore degenerates, and $i\circ j$ is regular by (1.7.2).

By (1.6.1), the sequence:
\[
 j^*(\underline{N}_{X/Y})\longrightarrow \underline{N}_{X/Z}\longrightarrow \underline{N}_{Y/Z}\longrightarrow0
\]
is exact. Since $i\circ j$ is regular, the three conormal sheaves are locally free of finite type; moreover, their rank at a point is the

%%%%%%%%%%  scan idx 437  |  folio 424  %%%%%%%%%%
\SGAPrintPage{424}%
number of elements in a minimum system of generators of the corresponding ideal; by the preceding proof, we have
\[
 \rang(\underline{N}_{X/Z})=\rang(\underline{N}_{X/Y})+\rang(\underline{N}_{Y/Z}).
\]
It follows that the kernel of $j^*(\underline{N}_{X/Y})\longrightarrow \underline{N}_{X/Z}$ is locally free of rank zero, and hence that the sequence (1.6.1) is exact.

ii) We may suppose $J/J^2$ free over $C$, and $I/I^2$ free over $B$; moreover, the hypothesis on the sequence of conormal sheaves allows us to suppose that $I/I^2\otimes_A A/J$ is a direct summand of $J/J^2$. We can therefore find elements $f_1,\ldots,f_n$ of $A$ lifting a basis of $I/I^2$, and elements $g_1,\ldots,g_m$ of $A$ such that $f_1,\ldots,f_n,g_1,\ldots,g_m$ lift a basis of $J/J^2$, the images of $g_1,\ldots,g_m$ generating $K/K^2$; after further localization, we may suppose that $f_1,\ldots,f_n$ form a minimum system of generators of $I$, that $f_1,\ldots,f_n,g_1,\ldots,g_m$ form a minimum system of generators of $J$, and that $\bar g_1,\ldots,\bar g_m$ generate $K$. We can then construct the diagram (1.7.1) and use the spectral sequence (1.7.3).

By (1.7.2), since $i$ is regular, we have $E^2_{pq}=0$ for $q\geq1$; we therefore obtain edge-isomorphisms: $E^2_{p,0}\simeq\Tor_p^{A_0}(C_0,A)$. Since $i\circ j$ is regular, this implies by (1.7.2) that $E^2_{p,0}=0$ for $p\geq1$, and hence that $j$ is regular, again by (1.7.2).

iii) This last property is proved in EGA IV 19.1.5 (iv), since in the noetherian case the definition in EGA IV coincides with the one we adopt here.

\sgaStableTarget{sga6.e07.prop.1-8}\textbf{Proposition 1.8.} \uline{Let} \(X\) \uline{be a scheme}, \(\underline{E}\) \uline{an} \(\underline{O}_X\)-\uline{Module locally free of finite type}, \(u:\underline{E}\to\underline{O}_X\) \uline{a regular homomorphism}, \(\underline{J}=u(\underline{E})\). \uline{Denote by} \(Y\)

%%%%%%%%%%  scan idx 438  |  folio 425  %%%%%%%%%%
\SGAPrintPage{425}%
\uline{the closed subscheme of} \(X\) \uline{defined by} \(\underline{J}\), \(X'\) \uline{the scheme obtained by blowing up} \(Y\), \(P\) \uline{the projective bundle associated with} \(\underline{E}\), \(p\) \uline{the structural morphism} \(P\to X\). \uline{On} \(P\), \uline{consider the canonical exact sequence}:
\sgaStableTarget{sga6.e07.eq.1-8-1}\[
0\to \underline{F}\to p^*(\underline{E})\to\underline{O}_P(1)\to0,
\]
\uline{and the composite morphism}:
\[
 v:\underline{F}\longrightarrow p^*(\underline{E})\xrightarrow{p^*(u)}\underline{O}_P .
\]

\uline{Then}:
\begin{enumerate}[label=\roman*)]
\item \uline{The formation of} \(X'\) \uline{commutes with every base change} \(f:T\to X\) \uline{such that} \(f^*(u)\) \uline{is regular}.
\item \uline{The closed immersion} \(j:X'\to P\) \uline{defined by the surjective homomorphism} \(\mathbf{S}(\underline{E})\longrightarrow\bigoplus_{n\ge0}\underline{J}^n\) \uline{defined by} \(u\) \uline{in degree} \(1\) \uline{is a regular immersion}, \uline{its ideal being the image of the regular morphism} \(v\); \uline{in particular, if} \(\operatorname{codim}(i)\) \uline{is constant}, \(\operatorname{codim}(j)=\operatorname{codim}(i)-1\).
\item \uline{The} \(\mathbf{S}(\underline{E})\)-\uline{Module} \(\bigoplus_{n\ge0}\underline{J}^n\) \uline{is perfect as a graded} \(\mathbf{S}(\underline{E})\)-\uline{Module}, \uline{i}.\uline{e}. \uline{locally there exists a finite}, \uline{graded resolution of} \(\bigoplus_{n\ge0}\underline{J}^n\) \uline{by graded} \(\mathbf{S}(\underline{E})\)-\uline{Modules locally free of finite type} (VI 1.2).
\end{enumerate}

i) Since \(X'=\Proj(\bigoplus_{n\ge0}\underline{J}^n)\), the result is an immediate consequence of the relations (1.2.3) and of EGA II.

ii) Since the stated properties are local on \(X\), we may suppose \(X\) affine, and moreover \(\underline{E}\) free over \(X\). Put \(X=\Spec(A)\).

First we give the proof in the standard case \(A=\mathbb Z[T_1,\ldots,T_d]\), with \(E\) endowed with a basis \(e_1,\ldots,e_d\) such that \(u(e_i)=T_i\). Then \(P=\Proj\bigl(\mathbb Z[T_1,\ldots,T_d,X_1,\ldots,X_d]\bigr)\), the grading being taken with respect to the \(X_i\). To show that \(v\) is a regular morphism whose image is the Ideal of \(j\), it suffices to do so locally on \(P\), hence on the

%%%%%%%%%%  scan idx 439  |  folio 426  %%%%%%%%%%
\SGAPrintPage{426}%
\(D^+(X_i)\); let us work, for example, on \(D^+(X_1)\). Then \(p^*(\underline{E})\) is defined by the module \(\mathbb Z[T_1,\ldots,T_d,X_2/X_1,\ldots,X_d/X_1]\otimes E\), and \(\underline{O}_P(1)\) by \(\mathbb Z[T_1,\ldots,T_d,X_2/X_1,\ldots,X_d/X_1]\), the canonical morphism sending \(e_i\) to \(X_i/X_1\) for every \(i\). Then \(F\) has as a basis the elements \(e_i-e_1.X_i/X_1\) for \(i>1\).

We have \(v(e_i-e_1.X_i/X_1)=T_i-T_1.X_i/X_1\); put \(T'_i=T_i-T_1X_i/X_1\).
Since the sequence of the \(T_i\) is regular in \(\mathbb Z[T_1,\ldots,T_d,X_2/X_1,\ldots,X_d/X_1]\), the same is true of the sequence \(T_1,T'_2,\ldots,T'_d\) deduced from it by an invertible matrix, and hence also of the subsequence \(T'_2,\ldots,T'_d\) (we are in a noetherian case). Consequently, \(v\) is a regular morphism.

We must therefore show that the \(T'_i\) generate the Ideal of \(j\). In \(\bigoplus_{n\ge0}J^n\), for \(m\le p\), we shall denote by \(x^{(m)}\) the element \(x\) of \(J^p\), considered as having degree \(m\). Over \(D^+(X_1)\), \(X'=\Spec\left(\left(\bigoplus_{n\ge0}J^n\right)_{(T_1^{(1)})}\right)\), and \(j\) is defined by the morphism sending \(X_i\) to \(T_i^{(1)}\). Under this morphism, the image of \(T'_i\) is \(T_i^{(0)}-T_1^{(0)}.T_i^{(1)}/T_1^{(1)}=\bigl((T_iT_1)^{(1)}-(T_1T_i)^{(1)}\bigr)/T_1^{(1)}=0\). Thus the \(T'_i\) belong to the Ideal of \(j\).

Since every element of \(\mathbb Z[T_1,\ldots,T_d,X_2/X_1,\ldots,X_d/X_1]\) can be expressed as a polynomial in \(T_1,T'_2,\ldots,T'_d,X_2/X_1,\ldots,X_d/X_1\), to show that the \(T'_i\) generate the Ideal of \(j\) it suffices to show that no nonzero element of \(\mathbb Z[T_1,X_2/X_1,\ldots,X_d/X_1]\) belongs to the Ideal of \(j\), which amounts to showing that the homomorphism
\[
\mathbb Z[T_1,X_1,\ldots,X_d]\longrightarrow\bigoplus_{n\ge0}J^n
\]
is injective, since \(T_1^{(1)}\) is not a zero divisor in \(\bigoplus_{n\ge0}J^n\). Grading \(\mathbb Z[T_1,X_1,\ldots,X_d]\) by assigning degree \(0\) to \(T_1\) and degree \(1\) to the \(X_i\), the kernel of the homomorphism under consideration is a graded ideal. Now it is clear that
